\clearpage\thispagestyle{empty}
\pdfbookmark[0]{Sources for Further Remarks on Easter (1807)}{osterfest1807-sources}
\begin{center}\textbf{Sources for Further Remarks on Easter (1807)}\end{center}
\noindent This AI translation follows Gauss’s \textit{Noch etwas über die Bestimmung des Osterfestes} in \textit{Werke}, vol.~VI (Göttingen, 1874), pp.~82--86. The first printing appeared in \textit{Braunschweigisches Magazin}, vol.~20, 12 September 1807, cols.~589--596, under the name Dr.~Gauss. \href{https://www.digitale-sammlungen.de/de/view/bsb10531610?page=303}{First printing}. \href{https://gdz.sub.uni-goettingen.de/id/PPN236007467}{Collected German text}.
\par\smallskip\noindent In the worked example for 1801, the first printing gives 90 and 7300 in the fourth-rule table (col.~592), where the collected text gives 48 and 7258 (p.~84). Both printings give 6 as the resulting fourth remainder. The collected text supplies the bracketed references to pp.~75, 78 and 80 of this volume. For a study of the successive versions of the Easter rules, see Reinhold Bien, \href{https://doi.org/10.1007/s00407-004-0078-5}{\textit{Gauß and Beyond: The Making of Easter Algorithms}}, \textit{Archive for History of Exact Sciences} 58 (2004), pp.~439--452, especially pp.~441--443.
\par\smallskip\noindent Translation and typesetting produced by AI.
\clearpage
